\section{Proofs and Operator Boundaries}
\label{sec:archive-proofs}

\subsection{Finite Attenuation of a Normalized Measure}

Fix a ray, its sample grid, and all primitive kernels.
Let $D_j>0$ be the total kernel mass of primitive $j$, and $N_j(b)\in[0,D_j]$ its
pre-boundary mass under the same linear interpolation as the deployed CDF.
For positive weights $w_j$, put
\begin{equation}
\begin{gathered}
D=\sum_jw_jD_j,\qquad N=\sum_jw_jN_j,\\
F=N/D,\qquad r_j=w_jD_j/D,\qquad C_j=N_j/D_j.
\end{gathered}
\end{equation}
Differentiating the quotient gives
\begin{equation}
\begin{aligned}
\frac{\partial F}{\partial\log w_j}
&=w_j\,\frac{N_jD-ND_j}{D^2}\\
&=\frac{w_jD_j}{D}\left(\frac{N_j}{D_j}-\frac ND\right)
=r_j(C_j-F).
\end{aligned}
\end{equation}
After $w_j$ is replaced by $vw_j$, $0<v<1$, the numerator and denominator become
$N_v=N-(1-v)w_jN_j$ and $D_v=D-(1-v)w_jD_j$.
Hence
\begin{equation}
F_v=\frac{F-(1-v)r_jC_j}{1-(1-v)r_j},
\end{equation}
and subtraction yields
\begin{equation}
F_v-F=\frac{(1-v)r_j(F-C_j)}{1-(1-v)r_j}.
\label{eq:archive-finite}
\end{equation}
The denominator is at least $v>0$.
Consequently, attenuation decreases $F$ iff $C_j>F$, increases it iff $C_j<F$,
and leaves it unchanged if $C_j=F$. At $v=1$ there is no change.
A zero-mass component has no effect.

For a family $\mathcal J$ uniformly scaled by $v$, replace $w_jD_j$ and $w_jN_j$ with
$\sum_{j\in\mathcal J}w_jD_j$ and $\sum_{j\in\mathcal J}w_jN_j$.
This proves the family identity with its own weighted conditional CDF.
It does not assert magnitude-independent signs for arbitrary heterogeneous multi-component changes.
Those changes can alter the balance between favorable and adverse components.

\paragraph{A two-component example.}
Take equal total contributions with $C_1=.8$ and $C_2=.2$.
The mixture has $F=.5$. Halving the early component gives $F_v=.4$; halving the later component
instead gives $F_v=.6$. Both operations reduce a positive weight.
Their opposite effects follow from which part of the normalized measure is reduced.

Uniformly duplicating every component leaves the normalized categorical distribution unchanged,
whereas duplicating only one family generally changes its relative contribution.
This is the remaining family-count dependence tested by M40/M41.
Global scale invariance should therefore not be interpreted as arbitrary discretization invariance.

\subsection{Median Decisions and Component Accounting}

For ordered samples, let $p_k>0$ and $F_k=\sum_{\ell\leq k}p_\ell$.
The implementation finds the first $k$ with $F_k\geq .5$ and linearly interpolates the crossing.
Within an interior positive-mass segment, the CDF is strictly increasing, so
\begin{equation}
\widehat d<b\quad\Longleftrightarrow\quad F(b)>.5.
\label{eq:archive-categorical-boundary}
\end{equation}
Boundary equality is not an early event.
At or outside the box endpoints, the implementation's endpoint convention must be applied explicitly;
a universal quantile identity for distributions with atoms or plateaus would need corresponding tie rules.

Interpolation is linear in sample mass.
Since $p_k=p_{k,\mathcal A}+p_{k,\mathrm{child}}$, it follows that
$F(b)=F_{\mathcal A}(b)+F_{\mathrm{child}}(b)$ using exactly the same interpolation coefficient.
M44 checks this accounting against the rendered early decision on all 99,208 valid source rays:
event agreement is exact and the maximum numerical decomposition residual is $2.98\times10^{-7}$.

\subsection{Why Expected Depth Does Not Control the Minimum}

Consider target depth $d^\star=1$ and early boundary $b=.8$.
A reference distribution puts all mass at $1+\delta$, with $\delta>0$.
Its expected-depth error is $\delta$ and its hard minimum is not early.
Construct a candidate with mass
\begin{equation}
\varepsilon=\frac{\delta}{.5+\delta}
\end{equation}
at depth $.5$ and the remaining mass at $1+\delta$.
Its expected depth equals one exactly:
$\varepsilon(.5)+(1-\varepsilon)(1+\delta)=1$.
Expected-depth error improves to zero, yet the hard minimum becomes $.5<b$.
As $\delta\to0$, the added early mass also tends to zero.
Thus arbitrarily small changes in expectation can accompany a discontinuous change in the minimum event.

This construction concerns a support-minimum operator, not the categorical median.
A guarantee for a literal minimum requires a constraint on early support.
A guarantee for a categorical median requires a margin on its ordered CDF.
A mean smooth-loss improvement supplies neither condition by itself.

\begin{table}[h]
\centering\small
\setlength{\tabcolsep}{4pt}
\begin{tabular}{p{.48\linewidth}rr}
\toprule
M51: M50 versus M8 & Same support & Deployment\\
\midrule
Added early rays (\%) & .594 & .659\\
Removed early rays (\%) & .104 & .114\\
Net early change (pp) & +.490 & +.545\\
Added rays with lower smooth error (\%) & 38.03 & 36.54\\
\bottomrule
\end{tabular}
\caption{Smooth/minimum diagnosis on 99,208 rays. Same support uses limited anchors plus children without
voxelization; deployment uses all anchors and the fixed $.06$\,m voxel realization.}
\end{table}

The mean smooth absolute error also changes from $.291706$ to $.292153$\,m,
so M50 exhibits both aggregate smooth regression and per-ray operator non-implication.
Smooth-error and same-support hard-depth changes correlate at $.038$.
Support realization changes 1.280\% of current-model early events, but the disagreement already exists
before that realization.

\subsection{Rigid Equivariance and Appearance Non-Interference}

Let $e_P$ depend only on canonical physical state $P$, and define $Q(P,T;x)=e_P(T^{-1}x)$.
For any rigid transform $g$,
\begin{equation}
Q(P,gT;gx)=e_P((gT)^{-1}gx)=e_P(T^{-1}x)=Q(P,T;x).
\end{equation}
For two canonical centers, rigid placement also gives
$\|Rc_j+t-(Rc_k+t)\|=\|c_j-c_k\|$ because $R^\top R=I$.
These statements require a common transform and unchanged canonical state.

Let visual parameters $\psi$ own $V_\psi$ and suppose the image optimizer receives only
$\operatorname{sg}(P)$ and updates $\psi$.
Since the physical query reads no component of $V_\psi$,
\begin{equation}
\frac{\partial Q(P,T;x)}{\partial\psi}=0.
\end{equation}
This is a structural property of the computation graph.
It remains meaningful even when the physical geometry, supplied pose, or visual reconstruction is inaccurate.

\subsection{Earlier Deletion and Reliability Identities}

For point sets $X'\subseteq X$, Eq.~\ref{eq:archive-literal} gives
$d_{X'}(r)\geq d_X(r)$, including the empty-set convention.
The early indicator can therefore only decrease under deletion.
Hit recall can either increase (an early occluder is removed) or decrease (the hit is removed).
This was the V7 motivation for evaluating coverage and return fidelity together.

For historical trajectory reliability, let
$C_\tau=\max_j |n_{\tau j}^{\top}e_j|/\max(d_{\tau j},\epsilon)$.
A uniform clearance change $\delta$ satisfies
\begin{equation}
|C_\tau(d+\delta)-C_\tau(d)|
\leq\max_j\frac{|n_{\tau j}^{\top}e_j|\,|\delta|}
{\max(d_j,\epsilon)\max(d_j+\delta,\epsilon)}.
\end{equation}
The result follows from the maximum operator and clipped denominator being 1-Lipschitz.
It is a deterministic sensitivity bound.

The earlier multi-horizon event $R(H,c)=
\Pr(\max_{t\leq H}C_t\leq c)$ is nondecreasing in budget $c$ and nonincreasing in prefix horizon $H$.
Convex interpolation between ordered budget knots preserves the first order; shared nonnegative interpolation
weights across ordered horizon surfaces preserve the second.
These structural orders are independent of probability calibration, which is evaluated on its own cohorts.
